%% kapa: exam=1 round=35 year=2024 period=2 form=A subject=appraisal_law subject_ko=감정평가관계법규 seq=23 number=23
%% source: Q-Net 감정평가사 기출문제 (cst00309) artlSeq=5213229 file=raw/1008/5213229/2교시 문제지 원본.pdf
%% conversion: kordoc 4.24.51 -> tools/gen1.py (auto); review=proofread-sonnet; flags=none
\begin{question}{23}
국유재산법령상 일반재산에 관한 설명으로 옳지 않은 것은?
\begin{choices}
  \item 국가가 매각한 일반재산을 일정기간 계속하여 점유ㆍ사용하는 경우에는 매각대금이 완납되기 전에 매각재산의 소유권을 이전할 수 있다.
  \item 일반재산을 매각한 경우에 매수자가 매각대금을 체납하면 그 매각계약을 해제할 수 있다.
  \item 일반재산의 매각대금이 3천만원을 초과하는 경우 매각대금을 5년 이내의 기간에 걸쳐 나누어 내게 할 수 있다.
  \item 일반재산을 용도를 지정하여 매각하는 경우에는 매수자는 매각일부터 10년 이상 지정된 용도로 활용하여야 한다.
  \item 부동산신탁을 취급하는 신탁업자에게 신탁하여 개발된 일반재산의 대부기간은 30년 이내로 할 수 있으며, 20년의 범위에서 한 차례만 연장할 수 있다.
\end{choices}
\end{question}
\answer{1}
